\documentclass[12pt,a4paper]{article}

\usepackage[utf8]{inputenc}
\usepackage[T1]{fontenc}
\usepackage[brazil]{babel}
\usepackage[a4paper,margin=2.5cm]{geometry}
\usepackage{setspace}
\usepackage{titlesec}
\usepackage{graphicx}
\usepackage{float}
\usepackage{enumitem}
\usepackage{array}
\usepackage{xcolor}
\usepackage{hyperref}
\usepackage{fancyhdr}
\usepackage{booktabs}

\hypersetup{
    colorlinks=true,
    linkcolor=blue,
    urlcolor=blue
}

\pagestyle{fancy}
\fancyhf{}
\fancyhead[L]{Relatório - Haiku OS}
\fancyhead[R]{Sistemas Operacionais}
\fancyfoot[C]{\thepage}

\begin{document}

%------------------------------------------------
% CAPA
%------------------------------------------------

\begin{titlepage}
    \centering
    
    {\Large \textbf{CENTRO UNIVERSITÁRIO MAURÍCIO DE NASSAU - UNINASSAU}}\\
    \vspace{2cm}
    
    {\Huge \textbf{HAIKU OS}}\\
    \vspace{1cm}
    
    {\Large Relatório da disciplina de Sistemas Operacionais}\\
    \vspace{2cm}
    
    \textbf{Alunos:}\\[0.5cm]
    
    André Gomes Moura - 01793413\\
    Carlos Vinícius Andrade Carvalho - 01581860\\
    João Pedro Machado da Silva - 01782106\\
    Jorge Vilela dos Santos Neto - 01778458\\
    
    \vfill
    
    Maceió - AL\\
    2026
    
\end{titlepage}

%------------------------------------------------
% SUMÁRIO
%------------------------------------------------

\tableofcontents
\newpage

%------------------------------------------------
% INTRODUÇÃO
%------------------------------------------------

\section{Introdução}

O Haiku OS é um sistema operacional de código aberto desenvolvido com foco em computação pessoal.
Inspirado diretamente no antigo sistema operacional BeOS, o projeto busca oferecer um ambiente moderno, leve, rápido e eficiente para computadores desktop.
O sistema prioriza simplicidade, desempenho e boa experiência do usuário, funcionando bem até mesmo em computadores mais modestos.
Além disso, o Haiku é mantido pela comunidade open source, permitindo colaboração de desenvolvedores do mundo inteiro.

%------------------------------------------------
% O QUE É O HAIKU
%------------------------------------------------

\section{O que é o Haiku OS}

O Haiku OS é um sistema operacional open source inspirado no BeOS, criado para oferecer rapidez, estabilidade e simplicidade.

\subsection{Sistema Operacional Open Source}

Por possuir código aberto, qualquer pessoa pode estudar, modificar e contribuir para o desenvolvimento do sistema.
O Haiku foi inspirado no \textbf{BeOS}, sistema operacional revolucionário lançado em 1995 pela empresa Be Inc., fundada por Jean-Louis Gassée, ex-executivo da Apple.

\subsection{Interface Moderna}

O sistema possui interface moderna, organizada e intuitiva, oferecendo facilidade de uso e boa experiência ao usuário.

\subsection{Leve, Rápido e Simples}

O Haiku consome poucos recursos do computador, possui inicialização rápida e responde rapidamente aos comandos do usuário.

\subsection{Foco em Desktop}

Seu foco principal é o uso em computadores pessoais, priorizando multitarefa, desempenho e facilidade de uso.

%------------------------------------------------
% HISTÓRIA
%------------------------------------------------

\section{História do Haiku}

O Haiku surgiu como um projeto open source criado para continuar o legado do antigo sistema operacional BeOS.
O BeOS foi desenvolvido na década de 1990 e ficou conhecido por ser extremamente rápido, leve e moderno para sua época.
Entre suas principais qualidades estavam:

\begin{itemize}
    \item Excelente multitarefa;
    \item Boa interface gráfica;
    \item Alto desempenho;
    \item Inicialização rápida;
    \item Ótima eficiência no uso do hardware.
\end{itemize}

Com o encerramento do projeto original, a comunidade decidiu criar o Haiku OS para recriar e modernizar o sistema, mantendo foco em desempenho, leveza e simplicidade.

%------------------------------------------------
% CARACTERÍSTICAS
%------------------------------------------------

\section{Características Principais}

\subsection{Arquitetura Híbrida e Modular}

O kernel do Haiku utiliza arquitetura híbrida, combinando características de kernels monolíticos e microkernels.

\subsection{Uso de C++ e Orientação a Objetos}

O Haiku faz forte uso da linguagem C++, diferente da maioria dos sistemas operacionais modernos.

\subsection{Baixa Latência (Soft Real-Time)}

O sistema possui foco em alta responsividade e baixa latência.

\subsection{Compatibilidade POSIX}

Apesar de independente, o Haiku possui compatibilidade com o padrão POSIX.

%------------------------------------------------
% ESTRUTURA
%------------------------------------------------

\section{Estrutura do Sistema Operacional}

O sistema operacional Haiku possui os seguintes componentes principais:

\begin{itemize}
    \item Kernel;
    \item Interface gráfica;
    \item Terminal;
    \item Gerenciador de arquivos;
    \item Drivers.
\end{itemize}

%------------------------------------------------
% KERNEL
%------------------------------------------------

\section{Kernel do Haiku}

O kernel é o núcleo do sistema operacional e funciona como o cérebro do computador.
Ele é responsável por controlar a comunicação entre software e hardware.

\subsection{Funções do Kernel}

\begin{itemize}
    \item Gerenciar memória;
    \item Controlar hardware;
    \item Executar programas;
    \item Controlar drivers e dispositivos.
\end{itemize}

%------------------------------------------------
% CARACTERÍSTICAS DO KERNEL
%------------------------------------------------

\section{Características do Kernel}

O kernel do Haiku apresenta diversas características importantes:

\begin{itemize}
    \item Leveza;
    \item Alto desempenho;
    \item Multitarefa eficiente;
    \item Gerenciamento de memória;
    \item Estabilidade;
    \item Otimização para desktop.
\end{itemize}

%------------------------------------------------
% LINGUAGENS
%------------------------------------------------

\section{Linguagens de Programação}

\subsection{C++}

Utilizada para alta performance e desenvolvimento de sistemas complexos.

\subsection{C}

Responsável pelo controle de baixo nível e interação direta com hardware.

\subsection{Python}

Usada para produtividade, prototipagem rápida e código limpo.

\subsection{Shell Script}

Utilizada para automação de tarefas e gerenciamento do sistema.

%------------------------------------------------
% VANTAGENS E DESVANTAGENS
%------------------------------------------------

\section{Vantagens e Desvantagens}

\subsection{Vantagens}

\begin{itemize}
    \item Sistema leve e rápido;
    \item Baixo consumo de recursos;
    \item Interface simples;
    \item Boa experiência em computadores mais fracos.
\end{itemize}

\subsection{Desvantagens}

\begin{itemize}
    \item Poucos programas disponíveis;
    \item Comunidade pequena;
    \item Compatibilidade limitada.
\end{itemize}

%------------------------------------------------
% CONCLUSÃO
%------------------------------------------------

\section{Conclusão}

O Haiku OS é um sistema operacional moderno, leve e eficiente, inspirado diretamente no antigo BeOS.
Seu foco principal está na simplicidade, desempenho e experiência do usuário no ambiente desktop.
Mesmo possuindo limitações como menor quantidade de programas e comunidade reduzida, o sistema se destaca por sua rapidez, estabilidade e otimização.
Dessa forma, o Haiku representa uma alternativa interessante para usuários que procuram um sistema operacional open source leve e funcional.

\end{document}